\subsection{Coxeter 群的刻画}

定理 1. (W,S) 是 Coxeter 系统，当且仅当它满足第 5 号的交换条件 (E)。

第 4 号的引理 3 表明，任意 Coxeter 系统都满足 (E)。

反之，假设 (E) 得到满足。设 G 为一个群，f 为从 S 到 G 的映射，使得 \((f(s)f(s'))^{m}=1\)，只要 s 和 \(s'\) 属于 S 且 \(ss'\) 的阶为有限数 m。由命题 5，存在从 W 到 G 的映射 g，使得

\[
g (w) = f (s _ {1}) \dots f (s _ {q})\tag{20}
\]

只要 \(w = s_{1} \ldots s_{q}\) 的长度为 q，就有上述等式。为了证明 (W,S) 是 Coxeter 系统，只需证明 g 是同态；这由公式

\[
g (s w) = f (s) g (w) \quad \text { for } s \in S, w \in W\tag{21}
\]

因为 S 生成 W。由第 5 号的命题 4，只可能有以下两种情形：

\begin{enumerate}
\def\labelenumi{\alph{enumi})}
\item
  \(l(sw) = l(w) + 1\)：若 \((s_1, \ldots, s_q)\) 是 \(w\) 的一个约化分解，则 \((s, s_1, \ldots, s_q)\) 是 \(sw\) 的一个约化分解，于是 (21)。
\item
  \(l(sw) = l(w) - 1\)：置 \(w' = sw\)；于是 \(w = sw'\) 且 \(l(sw') = l(w') + 1\)。由 a)，\(g(w) = f(s)g(sw)\)，因而 \(f(s)g(w) = g(sw)\)，因为 \((f(s))^{2} = 1\)。
\end{enumerate}

